\begin{frame}{이번 주차 정리}
  \begin{enumerate}
    \item \kb{16.1} — \textbf{두 층의 베이스라인}(무작위·GBDT)
      을 \textbf{둘 다} 넘어야 ``딥러닝이 필요했다''. 근거는
      \textbf{2축}(구조·규모) + \textbf{2숫자}(GBDT 비교, 학습
      곡선). digits 64차원 GBDT 0.944 $\succ$ CNN 0.939,
      784차원에서 \textbf{역전} 0.961 = 축 2(규모)의 증거.
    \item \kb{16.2} — 이번엔 \textbf{너가 심사위원}. 체크리스트
      + ``왜 DL인가'' 5번째 항목. \textbf{아키텍처-데이터
      매칭}이 3점 질문의 최빈 영역. ``test 0.95 vs 0.94는
      \textbf{잡음 수준}''을 $\mathbb{E}[\max]$($\approx0.05$) +
      표본 잡음(0.06)으로 \textbf{정량화}.
    \item \kb{16.3} — \textbf{채점 합산}(발표 50/보고서 30/
      리뷰 20)과 ``\textbf{리뷰 한 통이 94\%를 88\%로 바
      껐다}''(그룹 누수). Ch02~15를 \textbf{4단계 구조}
      (모델 정의$\to$손실$\to\theta$ 조정$\to$통제)와 개념
      지도 위에.
  \end{enumerate}
  \vskip0.8em
  \begin{block}{다음 학기 — ML2: 강화학습·로봇}
    정답 $y$가 있는 세계에서 벗어난다 — 정답 대신 \textbf{보상
    $r$}만 주어지고, \textbf{에이전트가 스스로 데이터를 만들어
    가는} 로봇 시뮬레이션의 세계. 4단계 구조는 \textbf{그대로} —
    바뀌는 것은 ①의 \textbf{모델}(정책 $\pi(a|s)$)과 ④의
    \textbf{데이터}(에이전트의 경험)뿐.
    \kq{같은 구조, 다른 질문.}
  \end{block}
\end{frame}
